\chapter{Variations of Energy}
\label{chap:dc-variations-of-energy}

Faithful transcription of do Carmo, \emph{Riemannian Geometry}. Each numbered
statement keeps do Carmo's number, kind and (name); proofs keep his explicit
cross-references ("by \cref{lem:dc-ch9-2-3}", "Theor. 3.7, Chap. 3") so the dependency graph
is recoverable from the text.

\section{Introduction}

In Chapter 3, geodesics were defined as curves with zero acceleration, and were
characterized by the fact that they locally minimize arc length. This chapter
presents a further characterization of a geodesic as a "solution of a variational
problem", adapting to Differential Geometry concepts and techniques from the
Calculus of Variations. The fundamental point of the chapter is the calculation
of the formula for the second variation of the energy of a geodesic (Section 2).
In Section 3 two geometric applications are made: the Theorem of Bonnet–Myers,
stating that a complete manifold whose curvature is positive and does not approach
zero is compact, with diameter estimated in terms of bounds on the curvature; and
an extension, due to A. Weinstein, of a theorem of Synge asserting the simple
connectivity of a compact, orientable manifold of even dimension whose curvature
is positive. Together with the theorem of Hadamard, these applications investigate
the influence of curvature on the topology of Riemannian manifolds, culminating in
the Sphere Theorem (Chap. 13).

\section{Formulas for the first and second variations of energy}

We make precise the idea of "neighboring curves" to a given curve.

\begin{definition}[variation]
  \label{def:dc-ch9-2-1}
  \dcref{ch9:2.1}
  Let $c:[0,a]\to M$ be a piecewise differentiable curve in a manifold $M$. A
  \emph{variation} of $c$ is a continuous mapping $f:(-\varepsilon,\varepsilon)\times[0,a]\to M$
  such that:
  - a) $f(0,t)=c(t)$, $t\in[0,a]$,
  - b) there exists a subdivision of $[0,a]$ by points
  $0=t_0<t_1<\cdots<t_{k+1}=a$, such that the restriction of $f$ to each
  $(-\varepsilon,\varepsilon)\times[t_i,t_{i+1}]$, $i=0,1,\dots,k$, is differentiable.
  
  A variation is said to be \emph{proper} if $f(s,0)=c(0)$ and $f(s,a)=c(a)$, for all
  $s\in(-\varepsilon,\varepsilon)$. If $f$ is differentiable, the variation is said to be
  \emph{differentiable}. For each $s\in(-\varepsilon,\varepsilon)$, the parametrized curve
  $f_s:[0,a]\to M$ given by $f_s(t)=f(s,t)$ is called a \emph{curve in the variation}, so
  a variation determines a family $f_s(t)$ of neighboring curves of $f_0(t)=c(t)$.
  The variation is proper exactly when the curves in this family have the same
  initial point $c(0)$ and endpoint $c(a)$.
  The parametrized differentiable curve given by $f_t(s)=f(s,t)$, $t$ fixed, is a
  \emph{transversal curve of the variation}. The velocity vector of a transversal curve
  at $s=0$, defined by $V(t)=\frac{\partial f}{\partial s}(0,t)$, is a (piecewise
  differentiable) vector field along $c(t)$ and is called the \emph{variational field} of
  $f$.
\end{definition}

\begin{proposition}[existence of a variation with given variational field]
  \label{prop:dc-ch9-2-2}
  \dcref{ch9:2.2}
  Given a piecewise differentiable field $V(t)$, along the piecewise differentiable
  curve $c:[0,a]\to M$, there exists a variation $f:(-\varepsilon,\varepsilon)\times[0,a]\to M$
  of $c$, such that $V(t)$ is the variational field of $f$; in addition, if
  $V(0)=V(a)=0$, it is possible to choose $f$ as a proper variation.
\end{proposition}
\begin{proof}
  Since $c([0,a])\subset M$ is compact it is possible to find a $\delta>0$
  such that $\exp_{c(t)}$ is well-defined for all $v\in T_{c(t)}M$ with $|v|<\delta$:
  for each $c(t)$ consider a totally normal neighborhood $W_t$ of $c(t)$ and the
  number $\delta_t>0$ associated to it (Theor. 3.7, Chap. 3); a finite subcover
  $W_1,\dots,W_n$ of $c([0,a])$ gives $\delta=\min(\delta_1,\dots,\delta_n)$.
  Consider $N=\max_{t\in[0,a]}|V(t)|$, $\varepsilon<\frac{\delta}{N}$ and define
  $f(s,t)=\exp_{c(t)}sV(t)$. Since $\exp_{c(t)}sV(t)=\gamma(1,c(t),sV(t))$ and the
  geodesic depends differentiably on the initial conditions, $f$ is piecewise
  differentiable, and $f(0,t)=c(t)$. The variational field of $f$ is
  
  $$
  \frac{\partial f}{\partial s}(0,t)=\frac{d}{ds}(\exp_{c(t)}sV(t))\Big|_{s=0}=(d\exp_{c(t)})_0 V(t)=V(t),
  $$
  
  and if $V(0)=V(a)=0$ then $f$ is proper. $\square$
  
  To compare the arc length of $c$ with that of neighboring curves we define
  $L:(-\varepsilon,\varepsilon)\to\mathbf{R}$ by
  
  $$
  L(s)=\int_0^a\Big|\frac{\partial f}{\partial t}(s,t)\Big|\,dt,
  $$
  
  the length of the curve $f_s(t)$. It is more convenient to work with the *energy
  function* $E(s)$ given by
  
  $$
  E(s)=\int_0^a\Big|\frac{\partial f}{\partial t}(s,t)\Big|^2\,dt,\qquad s\in(-\varepsilon,\varepsilon).
  $$
  
  For a curve $c:[0,a]\to M$ set
  
  $$
  L(c)=\int_0^a\Big|\frac{dc}{dt}\Big|\,dt,\qquad E(c)=\int_0^a\Big|\frac{dc}{dt}\Big|^2\,dt.
  $$
  
  Putting $f\equiv 1$ and $g=|\frac{dc}{dt}|$ in the Schwarz inequality
  $(\int_0^a fg\,dt)^2\le\int_0^a f^2\,dt\cdot\int_0^a g^2\,dt$, we obtain
  
  $$
  L(c)^2\le a\,E(c),
  $$
  
  with equality if and only if $g$ is constant, that is, if and only if $t$ is
  proportional to arc length.
\end{proof}

\begin{lemma}[minimizing geodesics minimize energy]
  \label{lem:dc-ch9-2-3}
  \dcref{ch9:2.3}
  \uses{cor:dc-ch3-3-9}
  Let $p,q\in M$ and let $\gamma:[0,a]\to M$ be a minimizing geodesic joining $p$ to
  $q$. Then, for all curves $c:[0,a]\to M$ joining $p$ to $q$,
  
  $$
  E(\gamma)\le E(c)
  $$
  
  with equality holding if and only if $c$ is a minimizing geodesic.
\end{lemma}
\begin{proof}
  From the considerations above,
  
  $$
  aE(\gamma)=(L(\gamma))^2\le(L(c))^2\le aE(c),
  $$
  
  which proves the first assertion. If equality holds, then $(L(c))^2=aE(c)$, so the
  parameter of $c$ is proportional to arc length, and $L(\gamma)=L(c)$, so $c$ is a
  minimizing geodesic (see \cref{cor:dc-ch3-3-9}). The converse is obvious.
  $\square$
\end{proof}

\begin{proposition}[formula for the first variation of the energy of a curve]
  \label{prop:dc-ch9-2-4}
  \dcref{ch9:2.4}
  Let $c:[0,a]\to M$ be a piecewise differentiable curve and let
  $f:(-\varepsilon,\varepsilon)\times[0,a]\to M$ be a variation of $c$. If
  $E:(-\varepsilon,\varepsilon)\to\mathbf{R}$ is the energy of $f$ then
  
  $$
  \frac{1}{2}E'(0)=-\int_0^a\Big\langle V(t),\frac{D}{dt}\frac{dc}{dt}\Big\rangle dt-\sum_{i=1}^k\Big\langle V(t_i),\frac{dc}{dt}(t_i^+)-\frac{dc}{dt}(t_i^-)\Big\rangle-\Big\langle V(0),\frac{dc}{dt}(0)\Big\rangle+\Big\langle V(a),\frac{dc}{dt}(a)\Big\rangle,\qquad(1)
  $$
  
  where $V(t)$ is the variational field of $f$, and
  $\frac{dc}{dt}(t_i^+)=\lim_{t\to t_i,\,t>t_i}\frac{dc}{dt}$,
  $\frac{dc}{dt}(t_i^-)=\lim_{t\to t_i,\,t<t_i}\frac{dc}{dt}$.
\end{proposition}
\begin{proof}
  By definition,
  $E(s)=\int_0^a\langle\frac{\partial f}{\partial t},\frac{\partial f}{\partial t}\rangle dt=\sum_{i=0}^k\int_{t_i}^{t_{i+1}}\langle\frac{\partial f}{\partial t},\frac{\partial f}{\partial t}\rangle dt$.
  Differentiating under the integral sign and using the symmetry of the Riemannian
  connection,
  
  $$
  \frac{d}{ds}\int_{t_i}^{t_{i+1}}\Big\langle\frac{\partial f}{\partial t},\frac{\partial f}{\partial t}\Big\rangle dt=2\Big\langle\frac{\partial f}{\partial s},\frac{\partial f}{\partial t}\Big\rangle\Big|_{t_i}^{t_{i+1}}-2\int_{t_i}^{t_{i+1}}\Big\langle\frac{\partial f}{\partial s},\frac{D}{dt}\frac{\partial f}{\partial t}\Big\rangle dt.
  $$
  
  Therefore,
  
  $$
  \frac{1}{2}\frac{dE}{ds}=\sum_{i=0}^k\Big\langle\frac{\partial f}{\partial s},\frac{\partial f}{\partial t}\Big\rangle\Big|_{t_i}^{t_{i+1}}-\int_0^a\Big\langle\frac{\partial f}{\partial s},\frac{D}{dt}\frac{\partial f}{\partial t}\Big\rangle dt.\qquad(2)
  $$
  
  Putting $s=0$ in (2) yields (1). $\square$
\end{proof}

\begin{proposition}[geodesics as critical points of energy]
  \label{prop:dc-ch9-2-5}
  \dcref{ch9:2.5}
  A piecewise differentiable curve $c:[0,a]\to M$ is a geodesic if and only if, for
  every proper variation $f$ of $c$, we have $\frac{dE}{ds}(0)=0$.
\end{proposition}
\begin{proof}
  If $c$ is a geodesic, $\frac{D}{dt}\frac{dc}{dt}=0$ and $c$ is regular;
  for $f$ proper, $V(0)=V(a)=0$, and all terms of (1) are zero. Conversely, suppose
  $\frac{dE}{ds}(0)=0$ for all proper variations. Let
  $V(t)=g(t)\frac{D}{dt}\frac{dc}{dt}$, where $g:[0,a]\to\mathbf{R}$ is piecewise
  differentiable with $g(t)>0$ if $t\ne t_i$ and $g(t_i)=0$, $i=0,\dots,k+1$.
  Constructing a variation
  with this $V(t)$,
  
  $$
  \frac{1}{2}\frac{dE}{ds}(0)=-\int_0^a g(t)\Big\langle\frac{D}{dt}\frac{dc}{dt},\frac{D}{dt}\frac{dc}{dt}\Big\rangle dt=0,
  $$
  
  so $\frac{D}{dt}\frac{dc}{dt}=0$ on each $(t_i,t_{i+1})$, that is, $c$ is a geodesic
  on each $(t_i,t_{i+1})$. To handle the points $t_i$, take another variational field
  $\overline V(t)$ with $\overline V(0)=\overline V(a)=0$ and, for $t\ne 0,a$,
  $\overline V(t_i)=\frac{dc}{dt}(t_i^+)-\frac{dc}{dt}(t_i^-)$. Then
  
  $$
  0=\frac{1}{2}\frac{dE}{ds}(0)=-\sum_{i=1}^k\Big|\frac{dc}{dt}(t_i^+)-\frac{dc}{dt}(t_i^-)\Big|^2,
  $$
  
  so $c$ is of class $C^1$ at each $t_i$. Since $\frac{D}{dt}\frac{dc}{dt}=0$ at
  $t_i$, $c$ satisfies the geodesic equation on $(0,a)$; by uniqueness of solutions
  to ordinary differential equations, $c\in C^\infty$ and is a geodesic. $\square$
\end{proof}

\begin{remark}[Remark 2.6]
  \label{rem:dc-ch9-2-6}
  \dcref{ch9:2.6}
  \uses{prop:dc-ch9-2-5}
  \cref{prop:dc-ch9-2-5} furnishes a characterization of geodesics as critical points of
  the energy for all proper variations. In this sense geodesics can be thought of as
  solutions to a variational problem; such a characterization is not local but
  involves the behavior of the geodesic as a whole.
\end{remark}

\begin{remark}[Remark 2.7]
  \label{rem:dc-ch9-2-7}
  \dcref{ch9:2.7}
  There is an analogy between the first variation of a proper variation and the usual
  derivative of a function on a differentiable manifold. It is natural to think of
  the set $\Omega_{p,q}$ of piecewise differentiable curves $c$ joining two points
  $p$ and $q$ of $M$ as a manifold, a tangent vector at the point $c$ being a
  piecewise differentiable vector field $V$ along $c$ which vanishes at the
  extremities of $c$. The energy $E$ is then a differentiable function on such a
  manifold, and $\frac{dE}{ds}(0)$ is the derivative of $E$ in the direction of $V$;
  geodesics joining $p$ to $q$ are critical points of $E$. The tangent space at a
  point $c$ does not have finite dimension, so local parametrizations by open subsets
  of $\mathbf{R}^n$ are impossible; this suggests manifolds with infinitely many
  dimensions (consult Palais, R., "Morse theory on Hilbert manifolds", Topology 2
  (1963), 299–340).
\end{remark}

\begin{proposition}[formula for the second variation]
  \label{prop:dc-ch9-2-8}
  \dcref{ch9:2.8}
  Let $\gamma:[0,a]\to M$ be a geodesic and let
  $f:(-\varepsilon,\varepsilon)\times[0,a]\to M$ be a proper variation of $\gamma$. Let $E$
  be the energy function of the variation. Then
  
  $$
  \frac{1}{2}E''(0)=-\int_0^a\Big\langle V(t),\frac{D^2 V}{dt^2}+R\Big(\frac{d\gamma}{dt},V\Big)\frac{d\gamma}{dt}\Big\rangle dt-\sum_{i=1}^k\Big\langle V(t_i),\frac{DV}{dt}(t_i^+)-\frac{DV}{dt}(t_i^-)\Big\rangle,\qquad(3)
  $$
  
  where $V$ is the variational field of $f$, $R$ is the curvature of $M$ and
  $\frac{DV}{dt}(t_i^+)=\lim_{t\to t_i,\,t>t_i}\frac{DV}{dt}$,
  $\frac{DV}{dt}(t_i^-)=\lim_{t\to t_i,\,t<t_i}\frac{DV}{dt}$.
\end{proposition}
\begin{proof}
  Taking the derivative of (2),
  
  $$
  \frac{1}{2}\frac{d^2 E}{ds^2}=\sum_{i=0}^k\Big\langle\frac{D}{ds}\frac{\partial f}{\partial s},\frac{\partial f}{\partial t}\Big\rangle\Big|_{t_i}^{t_{i+1}}+\sum_{i=0}^k\Big\langle\frac{\partial f}{\partial s},\frac{D}{ds}\frac{\partial f}{\partial t}\Big\rangle\Big|_{t_i}^{t_{i+1}}-\int_0^a\Big\langle\frac{D}{ds}\frac{\partial f}{\partial s},\frac{D}{dt}\frac{\partial f}{\partial t}\Big\rangle dt-\int_0^a\Big\langle\frac{\partial f}{\partial s},\frac{D}{ds}\frac{D}{dt}\frac{\partial f}{\partial t}\Big\rangle dt.
  $$
  
  Putting $s=0$, the first and third terms are zero, since $f$ is proper and
  $\gamma$ is a geodesic. Because
  
  $$
  \frac{D}{ds}\frac{D}{dt}\frac{\partial f}{\partial t}=\frac{D}{dt}\frac{D}{ds}\frac{\partial f}{\partial t}+R\Big(\frac{\partial f}{\partial t},\frac{\partial f}{\partial s}\Big)\frac{\partial f}{\partial t},
  $$
  
  we have, at $s=0$,
  $\frac{D}{ds}\frac{D}{dt}\frac{\partial f}{\partial t}=\frac{D^2}{dt^2}V+R(\frac{d\gamma}{dt},V)\frac{d\gamma}{dt}$.
  Further use of the fact that the variation is proper yields
  
  $$
  \sum_{i=0}^k\Big\langle\frac{\partial f}{\partial s},\frac{D}{ds}\frac{\partial f}{\partial t}\Big\rangle\Big|_{t_i}^{t_{i+1}}=-\sum_{i=1}^k\Big\langle V(t_i),\frac{DV}{dt}(t_i^+)-\frac{DV}{dt}(t_i^-)\Big\rangle.\qquad(4)
  $$
  
  Putting all these facts together, we obtain (3). $\square$
\end{proof}

\begin{remark}[Remark 2.9]
  \label{rem:dc-ch9-2-9}
  \dcref{ch9:2.9}
  \uses{prop:dc-ch9-2-8}
  If the variation is not proper, the first term at the beginning of \cref{prop:dc-ch9-2-8}
  need not be zero. Taking this with the corresponding terms at $i=0$ and $i=k+1$ of
  the second member of (4), we obtain
  
  $$
  \frac{1}{2}E''(0)=-\int_0^a\Big\langle V(t),\frac{D^2 V}{dt^2}+R\Big(\frac{d\gamma}{dt},V\Big)\frac{d\gamma}{dt}\Big\rangle dt-\sum_{i=1}^k\Big\langle V(t_i),\frac{DV}{dt}(t_i^+)-\frac{DV}{dt}(t_i^-)\Big\rangle
  $$
  
  $$
  -\Big\langle\frac{D}{ds}\frac{\partial f}{\partial s},\frac{d\gamma}{dt}\Big\rangle(0,0)+\Big\langle\frac{D}{ds}\frac{\partial f}{\partial s},\frac{d\gamma}{dt}\Big\rangle(0,a)-\Big\langle V(0),\frac{DV}{dt}(0)\Big\rangle+\Big\langle V(a),\frac{DV}{dt}(a)\Big\rangle.\qquad(5)
  $$
\end{remark}

\begin{remark}[the index form]
  \label{rem:dc-ch9-2-10}
  \dcref{ch9:2.10}
  It is often convenient to write (5) using, on each interval where $V$ is
  differentiable,
  $\frac{d}{dt}\langle V,\frac{DV}{dt}\rangle=\langle V,\frac{D^2 V}{dt^2}\rangle+\langle\frac{DV}{dt},\frac{DV}{dt}\rangle$.
  Taking a geodesic $\gamma:[0,a]\to M$ and a partition
  $0=t_0<t_1<\cdots<t_{k+1}=a$, this gives
  
  $$
  \frac{1}{2}E''(0)=\int_0^a\{\langle V',V'\rangle-\langle R(\gamma',V)\gamma',V\rangle\}dt-\Big\langle\frac{D}{ds}\frac{\partial f}{\partial s},\gamma'\Big\rangle(0,0)+\Big\langle\frac{D}{ds}\frac{\partial f}{\partial s},\gamma'\Big\rangle(0,a).
  $$
  
  For reasons which will be clear later, it is convenient to write
  
  $$
  \int_0^a\{\langle V',V'\rangle-\langle R(\gamma',V)\gamma',V\rangle\}dt=I_a(V,V).
  $$
  
  If the variation is proper, $I_a(V,V)=\frac{1}{2}E''(0)$ depends only on $V$. In
  the general case,
  
  $$
  \frac{1}{2}E''(0)=I_a(V,V)-\Big\langle\frac{D}{ds}\frac{\partial f}{\partial s},\gamma'\Big\rangle(0,0)+\Big\langle\frac{D}{ds}\frac{\partial f}{\partial s},\gamma'\Big\rangle(0,a).\qquad(6)
  $$
\end{remark}

\begin{remark}[Remark 2.11]
  \label{rem:dc-ch9-2-11}
  \dcref{ch9:2.11}
  It is possible to establish formulas for the first and second variations of the
  arc length function $L(s)$. The expressions are entirely analogous to those for the
  energy and, since they will not be used here, are not treated (see [dC 2]
  paragraph 5.4).
\end{remark}

\begin{remark}[Remark 2.12]
  \label{rem:dc-ch9-2-12}
  \dcref{ch9:2.12}
  \uses{rem:dc-ch9-2-7}
  The analogy mentioned in \cref{rem:dc-ch9-2-7} can be carried further, considering
  $2I_a(V,V)$ as the hessian $d^2 E(V,V)$ of the energy function
  $E:\Omega_{p,q}\to\mathbf{R}$ at the critical point $\gamma\in\Omega_{p,q}$ with
  respect to the "vector" $V$.
\end{remark}

\section{The theorems of Bonnet–Myers and of Synge–Weinstein}

We now go into some applications of the formula for the second variation of the
energy.

\begin{theorem}[Bonnet–Myers]
  \label{thm:dc-ch9-3-1}
  \dcref{ch9:3.1}
  \uses{lem:dc-ch9-2-3}
  Let $M^n$ be a complete Riemannian manifold. Suppose that the Ricci curvature of
  $M$ satisfies
  
  $$
  \mathrm{Ric}_p(v)\ge\frac{1}{r^2}>0,
  $$
  
  for all $p\in M$ and for all $v\in T_p(M)$. Then $M$ is compact and the diameter
  $\mathrm{diam}(M)\le\pi r$.
\end{theorem}
\begin{proof}
  Let $p$ and $q$ be any two points in $M$. Since $M$ is complete, there
  exists a minimizing geodesic $\gamma:[0,1]\to M$ joining $p$ to $q$. It is enough
  to show that the length $\ell(\gamma)\le\pi r$, because then $M$ is bounded and
  complete, therefore compact; and since $d(p,q)\le\pi r$ for any $p,q$, it follows
  that $\mathrm{diam}(M)\le\pi r$. Suppose, to the contrary, that
  $\ell(\gamma)=\ell>\pi r$. Consider parallel fields $e_1(t),\dots,e_{n-1}(t)$ along
  $\gamma$ which are orthonormal for each $t\in[0,1]$ and belong to the orthogonal
  complement of $\gamma'(t)$. Let $e_n(t)=\frac{\gamma'(t)}{\ell}$ and let $V_j$ be
  the vector field along $\gamma$ given by $V_j(t)=(\sin\pi t)e_j(t)$,
  $j=1,\dots,n-1$. Then $V_j(0)=V_j(1)=0$, so $V_j$ generates a proper variation of
  $\gamma$ with energy $E_j$. Using the formula for the second variation and the fact
  that $e_j$ is parallel,
  
  $$
  \frac{1}{2}E_j''(0)=-\int_0^1\langle V_j,V_j''+R(\gamma',V_j)\gamma'\rangle dt=\int_0^1\sin^2\pi t\,(\pi^2-\ell^2 K(e_n(t),e_j(t)))\,dt,
  $$
  
  where $K(e_n(t),e_j(t))$ is the sectional curvature at $\gamma(t)$ with respect to
  the plane generated by $e_n(t),e_j(t)$. Summing on $j$ and using the definition of
  Ricci curvature,
  
  $$
  \frac{1}{2}\sum_{j=1}^{n-1}E_j''(0)=\int_0^1\{\sin^2\pi t\,((n-1)\pi^2-(n-1)\ell^2\,\mathrm{Ric}_{\gamma(t)}(e_n(t)))\}\,dt.
  $$
  
  Since $\mathrm{Ric}_{\gamma(t)}(e_n(t))\ge\frac{1}{r^2}$ and $\ell>\pi r$, we have
  $(n-1)\ell^2\,\mathrm{Ric}_{\gamma(t)}(e_n(t))>(n-1)\pi^2$, hence
  
  $$
  \frac{1}{2}\sum_{j=1}^{n-1}E_j''(0)<\int_0^1\sin^2\pi t\,((n-1)\pi^2-(n-1)\pi^2)\,dt=0.
  $$
  
  As a result there exists an index $j$ such that $E_j''(0)<0$, which, by \cref{lem:dc-ch9-2-3},
  contradicts the fact that $\gamma$ is minimizing. Therefore $\ell\le\pi r$.
  $\square$
  
  In the corollaries that follow, we use some facts on the fundamental group and on
  covering spaces (see Massey [Ma], Chapters 2 and 5).
\end{proof}

\begin{corollary}[finiteness of the fundamental group]
  \label{cor:dc-ch9-3-2}
  \dcref{ch9:3.2}
  Let $M$ be a complete Riemannian manifold with $\mathrm{Ric}_p(v)\ge\delta>0$, for
  all $p\in M$ and all $v\in T_p(M)$. Then the universal cover of $M$ is compact. In
  particular, the fundamental group $\pi_1(M)$ is finite.
\end{corollary}
\begin{proof}
  Introducing on the universal cover $\pi:\tilde M\to M$ the covering metric
  (the metric such that $\pi$ is a local isometry), $\tilde M$ is complete and its
  Ricci curvature satisfies $\widetilde{\mathrm{Ric}}_p\ge\delta>0$. From the theorem,
  $\tilde M$ is compact. Hence the number of sheets of the covering is finite; since
  that is the number of elements in $\pi_1(M)$, we conclude that $\pi_1(M)$ is
  finite. $\square$
\end{proof}

\begin{corollary}[positive sectional curvature bound]
  \label{cor:dc-ch9-3-3}
  \dcref{ch9:3.3}
  Let $M$ be a complete Riemannian manifold with sectional curvature
  $K\ge\frac{1}{r^2}>0$. Then $M$ is compact, $\mathrm{diam}(M)\le\pi r$ and
  $\pi_1(M)$ is finite.
\end{corollary}

\begin{remark}[Remark 3.4]
  \label{rem:dc-ch9-3-4}
  \dcref{ch9:3.4}
  The hypothesis $K\ge\delta>0$ cannot be relaxed to $K>0$. Indeed, the paraboloid
  $\{(x,y,z)\in\mathbf{R}^3;z=x^2+y^2\}$ has curvature $K>0$, but is complete and
  non-compact.
\end{remark}

\begin{remark}[Remark 3.5]
  \label{rem:dc-ch9-3-5}
  \dcref{ch9:3.5}
  In fact it is not necessary that $K$ be bounded away from zero but only that $K$
  not approach zero too fast. In this respect, see E. Calabi, "On Ricci curvatures
  and geodesics", Duke Math. J. 34 (1967), 667–676 and R. Schneider, "Konvexe
  Flächen mit langsam abnehmender Krümmung", Archiv der Math. 23 (1972), 650–654.
\end{remark}

\begin{remark}[Remark 3.6]
  \label{rem:dc-ch9-3-6}
  \dcref{ch9:3.6}
  \uses{thm:dc-ch9-3-1}
  The estimate for the diameter given by \cref{thm:dc-ch9-3-1} cannot be improved. Indeed, the
  unit sphere $S^n\subset\mathbf{R}^{n+1}$ has constant sectional curvature equal to
  1 (therefore its Ricci curvature also is a constant equal to 1) and
  $\mathrm{diam}(S^n)=\pi$. This example is unique in the following sense: let $M^n$
  be complete with $\mathrm{Ric}_p(v)\ge 1/r^2$, for all $p\in M$ and all $v\in T_pM$;
  if $\mathrm{diam}(M)=\pi r$, then $M^n$ is isometric to the sphere $S^n$ of
  curvature $1/r^2$ (see S.Y. Cheng, "Eigenvalues comparison theorems and its
  geometric applications", Math. Z. 143 (1975), 289–297 and, for another proof, K.
  Shiohama, "A sphere theorem for manifolds of positive Ricci curvature", Trans.
  A.M.S. 275 (1983), 811–819).
\end{remark}

\begin{theorem}[Synge–Weinstein]
  \label{thm:dc-ch9-3-7}
  \dcref{ch9:3.7}
  Let $f$ be an isometry of a compact oriented Riemannian manifold $M^n$. Suppose
  that $M$ has positive sectional curvature and that $f$ preserves the orientation of
  $M$ if $n$ is even, and reverses it if $n$ is odd. Then $f$ has a fixed point, i.e.,
  there exists $p\in M$ with $f(p)=p$.
  
  In the proof of Theorem 3.7 we need the following fact from linear algebra.
\end{theorem}

\begin{lemma}[invariant vector of an orthogonal map]
  \label{lem:dc-ch9-3-8}
  \dcref{ch9:3.8}
  \uses{rem:dc-ch9-2-9, thm:dc-ch9-3-7}
  Let $A$ be an orthogonal linear transformation of $\mathbf{R}^{n-1}$ and suppose
  that $\det A=(-1)^n$. Then $A$ leaves invariant some non-zero vector of
  $\mathbf{R}^{n-1}$.
\end{lemma}
\begin{proof}
  If $n$ is even, $\det(A-\lambda I)$ is a real polynomial in $\lambda$ of
  odd degree $(=n-1)$, so $A$ has a real eigenvalue. Since $A$ is orthogonal, such
  eigenvalues are of the form $\pm 1$. Since the product of the complex eigenvalues
  of $A$ is non-negative, and $\det A=1$, at least one of the eigenvalues equals 1.
  If $n$ is odd, $\det A=-1$. Because the product of the complex eigenvalues is
  non-negative, there is at least one pair of real eigenvalues, one of which is
  positive, hence equal to 1. $\square$
  
  \emph{Proof of \cref{thm:dc-ch9-3-7}.} Suppose, to the contrary, that $f(q)\ne q$ for all $q\in M$.
  Let $p\in M$ such that $d(p,f(p))$ attains a minimum. Since $M$ is complete, there
  exists a normalized minimizing geodesic $\gamma:[0,\ell]\to M$, joining $p$ to
  $f(p)$, that is, $\gamma(0)=p$, $\gamma(\ell)=f(p)$. Let
  $\tilde A=P\circ df_p:T_p(M)\to T_p(M)$, where $P$ is the parallel transport along
  $\gamma$ from $f(p)=\gamma(\ell)$ to $p$. Then $\tilde A$ is an isometry. We show
  that $(f\circ\gamma)'(0)=\gamma'(\ell)$. Consider the geodesic $f\circ\gamma$ which
  joins $f(p)$ to $f^2(p)$. Let $p'=\gamma(t')$, $t'\ne 0$, $t'\ne\ell$, and
  $f(p')=f\circ\gamma(t')$. Since $f$ is an isometry, $d(p,p')=d(f(p),f(p'))$ and,
  from the triangle inequality,
  
  $$
  d(p',f(p'))\le d(p',f(p))+d(f(p),f(p'))=d(p',f(p))+d(p,p')=d(p,f(p)).
  $$
  
  Because $d(p,f(p))$ is a minimum, $d(p',f(p'))=d(p',f(p))+d(f(p),f(p'))$. Therefore
  the curve formed by $\gamma$ and $f\circ\gamma$ is a geodesic, hence
  $(f\circ\gamma)'(0)=\gamma'(\ell)$, as asserted. It follows that $\tilde A$ leaves
  $\gamma'(0)$ fixed, since
  
  $$
  \tilde A(\gamma'(0))=P\circ df_p(\gamma'(0))=P((f\circ\gamma)'(0))=P(\gamma'(\ell))=\gamma'(0).
  $$
  
  Let $A$ be the restriction of $\tilde A$ to the orthogonal complement of
  $\gamma'(0)$. Then $A$ is orthogonal on $\mathbf{R}^{n-1}$ and, since $P$ is an
  orientation-preserving isometry,
  
  $$
  \det A=\det\tilde A=\det(P\circ df_p)=(-1)^n,
  $$
  
  where in the last equality we use the hypothesis on $f$ and the fact that $P$
  preserves orientation. From the lemma, $A$ leaves a vector invariant. Let $e_1(t)$
  be a unit parallel field along $\gamma$ such that, for each $t$, $e_1(t)$ belongs to
  the orthogonal complement of $\gamma'(t)$ and $e_1(0)$ is invariant by $A$. Let
  $\beta(s)$, $s\in(-\varepsilon,\varepsilon)$, be a geodesic such that $\beta(0)=p$ and
  $\beta'(0)=e_1(0)$. Because $P\circ df_p(e_1(0))=e_1(0)$, we have
  $df_p(e_1(0))=e_1(\ell)$, that is, the geodesic $f\circ\beta$ is such that
  $f\circ\beta(0)=f(p)$ and $(f\circ\beta)'(0)=e_1(\ell)$. Let $h$ be a variation of
  $\gamma$ given by $h(s,t)=\exp_{\gamma(t)}(se_1(t))$,
  $s\in(-\varepsilon,\varepsilon)$, $t\in[0,\ell]$. Since $h(s,0)=\beta(s)$, then
  $h(s,\ell)=\exp_{f(p)}(se_1(\ell))=(f\circ\beta)(s)$. Therefore,
  
  $$
  V(t)=\frac{\partial}{\partial s}\exp_{\gamma(t)}(se_1(t))\Big|_{s=0}=e_1(t),
  $$
  
  hence $\frac{D^2 V}{dt^2}=0$. Using the formula for the second variation (see
  \cref{rem:dc-ch9-2-9}) and the fact that $\frac{\partial h}{\partial s}(0,t)=e_1(t)$,
  
  $$
  \frac{1}{2}\frac{d^2 E}{ds^2}(0)=-\int_0^\ell\Big\langle V(t),R\Big(\frac{d\gamma}{dt},V\Big)\frac{d\gamma}{dt}\Big\rangle dt+\Big\langle\frac{d\gamma}{dt},\frac{D}{ds}\frac{\partial h}{\partial s}\Big\rangle(0,\ell)-\Big\langle\frac{d\gamma}{dt},\frac{D}{ds}\frac{\partial h}{\partial s}\Big\rangle(0,0)=-\int_0^\ell K\Big(e_1(t),\frac{d\gamma}{dt}\Big)dt.
  $$
  
  Because the sectional curvature is positive, $\frac{1}{2}\frac{d^2 E}{ds^2}(0)<0$,
  and therefore $\frac{dE}{ds}$ is strictly decreasing on a neighborhood of zero.
  This shows that there exists a curve $c$ in the variation such that
  $(\ell(c))^2\le\ell E(c)<\ell E(\gamma)=\ell(\gamma)^2$. Since the curves in the
  variation join $q$ to $f(q)$, we obtain a contradiction to the fact that
  $d(p,f(p))$ is a minimum. $\square$
\end{proof}

\begin{remark}[Remark 3.9]
  \label{rem:dc-ch9-3-9}
  \dcref{ch9:3.9}
  The theorem remains true under the weaker hypothesis that $f$ is a conformal
  diffeomorphism. It is not known whether the theorem is still true if $f$ is merely
  a diffeomorphism. This would imply that $S^2\times S^2$ does not carry a metric of
  positive curvature, since the map $f$ which is the antipodal map on each factor
  preserves the orientation (each factor reverses the orientation) and does not have
  a fixed point. For more details, see A. Weinstein [We 2].
\end{remark}

\begin{corollary}[Synge]
  \label{cor:dc-ch9-3-10}
  \dcref{ch9:3.10}
  \uses{thm:dc-ch9-3-7, cor:dc-ch9-3-3}
  Let $M^n$ be a compact manifold with positive sectional curvature.
  - a) If $M^n$ is orientable and $n$ is even, then $M$ is simply connected.
  - b) If $n$ is odd, then $M^n$ is orientable.
\end{corollary}
\begin{proof}
  \textbf{a)} Let $\pi:\tilde M\to M$ be the universal covering of $M$. Introduce on
  $\tilde M$ the covering metric, and orient $\tilde M$ so that $\pi$ preserves
  orientation. Because $M$ is compact with positive curvature, $K\ge\delta>0$; since
  $\pi$ is a local isometry, the same curvature condition holds on $\tilde M$. Since
  $\tilde M$ is complete, by \cref{cor:dc-ch9-3-3}, $\tilde M$ is compact. Let
  $k:\tilde M\to\tilde M$ be a covering transformation, that is, $\pi\circ k=\pi$
  (see Massey [Ma], p. 159). Then $k$ is an isometry of $\tilde M$ which preserves
  orientation. Because $n$ is even, we can use the theorem to conclude that $k$ has a
  fixed point. But a covering transformation with a fixed point is the identity. It
  follows that the group of covering transformations of $\tilde M$ (isomorphic to the
  fundamental group of $M$; see Massey [Ma], p. 163) reduces to the identity.
  Therefore $M$ is simply connected.
  
  \textbf{b)} Suppose that $M$ is not orientable, and consider the orientable double cover
  $\overline M$ of $M$ (see Exercise 12 of Chap. 0). Introduce on $\overline M$ the
  covering metric. Since $\overline M$ is the double cover of a compact manifold,
  $\overline M$ is compact. Let $k$ be a covering transformation of $\overline M$,
  $k\ne\mathrm{id}$. Because $M$ is not orientable, $k$ is an isometry which reverses
  the orientation of $\overline M$. Since $n$ is odd, we can apply \cref{thm:dc-ch9-3-7} which
  guarantees that $k$ has a fixed point. Therefore $k=\mathrm{id}$, which is a
  contradiction. $\square$
\end{proof}

\begin{remark}[Remark 3.11]
  \label{rem:dc-ch9-3-11}
  \dcref{ch9:3.11}
  The real projective space $P^2(\mathbf{R})$ of dimension two, which is compact,
  non-orientable and not simply connected, is an example which shows the necessity of
  orientability in a) and odd dimension in b). On the other hand, $P^3(\mathbf{R})$
  which is compact, orientable and not simply connected is an example which shows the
  necessity of even dimension in (a).
\end{remark}
